\documentclass[10pt,a4paper]{amsart}
\usepackage[margin=19mm]{geometry}
\usepackage{amsmath,amssymb,mathtools}
\usepackage[scr=rsfs,cal=euler]{mathalfa}
\usepackage{xfrac}
\usepackage[unicode,hidelinks,bookmarks=false]{hyperref}
\newcommand{\ie}{i.e.\ }
\newcommand{\cf}{cf.\ }
\newcommand{\resp}{respectively\ }
\newcommand{\Subsection}{\subsection}
\providecommand{\nobreakdash}{\nobreak\hbox{-}}
\setlength{\emergencystretch}{2em}
\multlinegap0pt
\theoremstyle{plain}
\newtheorem{thm}{Theorem}
\newtheorem{lem}[thm]{Lemma}
\newtheorem{prop}[thm]{Proposition}
\newtheorem{cor}[thm]{Corollary}
\theoremstyle{definition}
\newtheorem{defn}[thm]{Definition}
\theoremstyle{remark}
\newtheorem{rem}[thm]{Remark}
\def\N{\mathbb{N}}
\def\Z{\mathbb{Z}}
\def\Q{\mathbb{Q}}
\def\R{\mathbb{R}}
\def\Id{\operatorname{Id}}
\def\LCM{\operatorname{LCM}}
\def\Im{\operatorname{Im}}
\title[Synchronization points: growth, asymptotics, congruences, and the synchronization zeta function]{Synchronization points: growth, asymptotics, congruences, and the synchronization zeta function: the form of the synchronization zeta function for two permutations of a finite set (printed as Theorem 30)}
\author{Alexander Fel'shtyn, Mateusz Slomiany}
\date{Source: arXiv:2601.21888v2 (CC BY 4.0) (CC BY 4.0)}
\begin{document}
\maketitle

\noindent\emph{Source and licence.} Synchronization points: growth, asymptotics, congruences, and the synchronization zeta function. Version arXiv:2601.21888v2 (CC BY 4.0), distributed under the Creative Commons Attribution 4.0 International licence, http://creativecommons.org/licenses/by/4.0/, as recorded in the arXiv licence metadata for this exact version and shown on the saved abstract page https://arxiv.org/abs/2601.21888v2. Full licence notice: licenses/arxiv-2601.21888-CC-BY-4.0.txt.\par\smallskip

\noindent\textit{Editorial scope.} Counted unit: the explicit form of the synchronization zeta function of two permutations of a finite set, printed in the article as Theorem 30 (source0.tex lines 1361--1375, one \texttt{thm} environment with source label \texttt{thm:synchronization-zeta-form-bijections-finite}), delivered together with the article's complete original proof of it (source0.tex lines 1376--1466, one \texttt{proof} environment of 91 lines immediately adjacent to the statement, with no gap and no deferral). The proof is the article's own argument: it evaluates the zeta function through the periodic orbit structure forced by the finite orders of the two permutations, groups the orbits by the roots of unity that occur, and reads off the exponents of the linear and quadratic factors, including the separate treatment of the case in which the common period is even. No mathematics is written, repaired or re-derived: statement and proof are the article's own text line for line, in the article's own notation, and no line of either is omitted, substituted or re-flowed.

Required context retained uncounted: the article's theorem on the periodicity of the sequence of synchronization numbers (source0.tex lines 1296--1307) and its theorem on the rationality of the logarithmic derivative of the synchronization zeta function with its labelled display (source0.tex lines 1327--1335). Those are the two results the delivered proof invokes, and nothing else from the article's development is used.

References. No citation command occurs in any retained line, so this package carries no bibliography block and no bibliography entry.

Editorial formatting only, in addition: an AMS \texttt{amsart} document that restates verbatim the article's own macro definitions needed by the retained lines and the article's own theorem-environment declarations with the same shared counter. Consequently the retained environments print here in the order in which they occur, so the counted theorem prints as Theorem 3 whereas the article prints it as Theorem 30; the two retained context theorems print here as Theorem 1 and Theorem 2, which the article prints as Theorems 28 and 29, so the relative order of the three delivered theorems is the same in both, the equations of the retained spans are numbered anew in order of appearance, and every cross-reference inside the retained text resolves to this wrapper's numbering. That is the only change to the numbering. No figure, no image-inclusion line and no drawing code occur in any retained line, so no artwork is delivered or needed.

Not retained: the article's introduction, its other theorems and examples, its dynamical and topological applications and its bibliography; none of that material is used as a step of the delivered proof, and the delivered proof is complete as the author wrote it. The package ships the article's concatenated source verbatim as \texttt{sources/193620/source0.tex}, and every delivered excerpt is an exact line range of that file. No mathematical statement, hypothesis, estimate or conclusion has been rewritten, repaired or added, and printed slips, if any, are preserved literally.
\begin{thm} \label{thm:synch-numbers-permutations-finite-set}
  Let $\sigma_i, \sigma_j : X \to X$ be two bijections, i.e. permutations of a finite set.
  Then the following hold
  \begin{enumerate}
  \item[\rm (i)] The sequence of synchronization numbers $\{\#S(\sigma_i^n, \sigma_j^n)\}_{n=1}^\infty$ is periodic.
  \item[\rm (ii)] If $L$ is the length of the period of the sequence of synchronization numbers
      then for $k=1,\ldots,\lfloor L/2 \rfloor$
      \begin{equation*} %\label{eq:S^k=S^L-k}
        \#S(\sigma_i^k, \sigma_j^k) = \#S(\sigma_i^{L-k}, \sigma_j^{L-k})
      \end{equation*}
  \end{enumerate}
\end{thm}

\begin{thm} \label{thm:sychronization-zeta-log-derivative-rational}
  Let $\sigma_i,\sigma_j: X \to X$ be two bijections, i.e. permutations, of a finite set.
  Then the logarithmic derivative of the synchronization zeta function $S_{\sigma_i, \sigma_j}(z)$
  is rational and is equal
  \begin{align} \label{eq:synchronization-zeta-log-derivative}
    \left[ \ln S_{\sigma_i, \sigma_j}(z) \right]' = \frac{S_{\sigma_i, \sigma_j}'(z)}{S_{\sigma_i, \sigma_j}(z)}
      = \frac{\sum_{n=1}^{L} \#S(\sigma_i^n, \sigma_j^n) z^{n-1}}{1-z^L}.
  \end{align}
\end{thm}

\begin{thm} \label{thm:synchronization-zeta-form-bijections-finite}
  Let $X$ be a finite set and $\sigma_1, \sigma_2: X \to X$ two bijections, i.e. permutations.
  Let $k_1, k_2 \in \N$ be such that $\sigma_1^{k_1} = \Id = \sigma_2^{k_2}$
  and denote $L := \LCM(k_1,k_2)$.
  Let $\omega := \exp(2\pi\iota/L)$ be the $L$-th root of $1$ and $\Re(\omega)$ its real part.

  Then the synchronization zeta function is of the form
  \begin{align} \label{eq:synchronization-permutations}
    S_{\sigma_1,\sigma_2}(z) = \begin{cases}
      (1-z)^{A_0} \cdot \prod_{k=1}^{\lfloor L/2 \rfloor} \left( z^2 - 2\Re(\omega^k)z+1 \right)^{A_k}, & 2 \!\not|\, L \\[1ex]
      (1-z)^{A_0}(1+z)^{A_{L/2}} \cdot \prod_{k=1}^{L/2-1} \left( z^2 - 2\Re(\omega^k)z+1 \right)^{A_k}, & 2 \,\vert\, L
    \end{cases},
  \end{align}
  where all $A_k \in \R$.
\end{thm}

\begin{proof}
  When we consider the logarithmic derivative of the synchronization zeta function formally, we can write
  \begin{align*}
    S_{\sigma_i, \sigma_j}(z) &= \exp \left( \int \frac{S_{\sigma_i, \sigma_j}'(z)}{S_{\sigma_i, \sigma_j}(z)} dz \right).
  \end{align*}
  We will integrate the logarithmic derivative using the form of Theorem~\ref{thm:sychronization-zeta-log-derivative-rational}.
  For simplicity, denote the numerator $\sum_{n=1}^{L} \#S(\sigma_i^n, \sigma_j^n) z^{n-1} =: P(z) \in \Z[z]$.
  The denominator has $L$ different roots which are $L$-th roots of $1$, i.e. $\omega^k, k=0,1,\ldots,L-1$
  therefore we factorize
  \begin{align*}
    z^L-1 = (z-1)(z-\omega)\cdots(z-\omega^{L-1}).
  \end{align*}
  Performing partial fraction decomposition we get
  \begin{align*}
    \frac{-P(z)}{z^L-1} &= \sum_{k=0}^{L-1} \frac{C_k}{z-\omega^k}
  \end{align*}
  and, after reducing the right side to a common denominator and denoting by $\widehat{(z-\omega^k)}$
  the missing term in the product,
  \begin{align*}
    \sum_{k=0}^{L-1} \frac{C_k}{z-\omega^k}
    &= \frac{\sum_{k=0}^{L-1} C_k \cdot (z-\omega^0) \cdots \widehat{(z-\omega^k)} \cdots (z-\omega^{L-1})}{z^L-1}.
  \end{align*}
  Observe that the derivative
  \begin{align*}
    (z^L-1)' = \left[ \prod_{k=0}^{L-1} (z-\omega^k) \right]'
      = \sum_{k=0}^{L-1} (z-\omega^0) \cdots \widehat{(z-\omega^k)} \cdots (z-\omega^{L-1})
  \end{align*}
  therefore, comparing numerators and substituting the roots one by one, for $k=0,\ldots,L-1$ we have
  \begin{align*}
    -P(\omega^k) &= C_k \cdot (z^L-1)'(\omega^k) = C_k \cdot L\omega^{k(L-1)} \\
    C_k &= \frac{-P(\omega^k)}{L\omega^{k(L-1)}}
  \end{align*}

  We are going to show that all $C_k$ are in fact real.
  If $L$ is even, then $\omega^{L/2} = \exp((2\pi\iota/L) \cdot (L/2)) = -1$
  is a single real root which we will consider separately.
  Regardless of parity of $L$, $1$ always is a single real root.
  We therefore have
  \begin{align*}
    C_0 = \frac{-P(1)}{L} \in \Q, \quad C_{L/2} = \frac{-P(-1)}{\pm L} \in \Q \text{ (for $L$ even)}.
  \end{align*}
  Now consider $L$ odd and $k$ such that $\omega^k \notin \R$. The imaginary part
  \begin{align*}
    \Im C_k &= \Im \left( -\frac{1}{L} \sum_{n=1}^L \#S(\sigma_i^n,\sigma_j^n) \omega^{k(n-1)+k(1-L)} \right) \\
%    &= -\frac{1}{L} \sum_{n=1}^L \#S(\sigma_i^n,\sigma_j^n) \Im \omega^{kn-kL} \\
%    &= -\frac{1}{L} \sum_{n=1}^L \#S(\sigma_i^n,\sigma_j^n) \Im\frac{\omega^{kn}}{\left(\omega^{L}\right)^k} \\
    &= -\frac{1}{L} \sum_{n=1}^L \#S(\sigma_i^n,\sigma_j^n) \Im \left(\omega^{k}\right)^n \\
  \intertext{Note that $(\omega^k)^L=1$ and for $n=1,\ldots,\lfloor L/2 \rfloor$ the terms appears in complex
      conjugate pairs $\Im(\omega^{k})^n = -\Im(\omega^{k})^{L-n}$ so that we can rewrite the last sum}
    &= -\frac{1}{L} \sum_{n=1}^{\lfloor L/2 \rfloor} \left[\#S(\sigma_i^n,\sigma_j^n)
      \Im \left(\omega^{k}\right)^n - \#S(\sigma_i^{L-n},\sigma_j^{L-n}) \Im \left(\omega^{k}\right)^n \right] \\
  \intertext{Using Theorem~\ref{thm:synch-numbers-permutations-finite-set} (ii) we get}
    &= -\frac{1}{L} \sum_{n=1}^{\lfloor L/2 \rfloor} \#S(\sigma_i^n,\sigma_j^n)
      \left[ \Im\left(\omega^k\right)^n - \Im\left(\omega^k\right)^n \right] \\
    &= 0
  \end{align*}
  therefore $C_k \in \R$ is in fact real.

  Because $z^L-1$ is a polynomial with integer coefficients, all its complex roots appear in pairs
  with their complex conjugates $\overline{\omega^k} = \omega^{-k} = \omega^{L-k}$.
  Grouping conjugate complex roots we get
  \begin{align*}
    \frac{-P(z)}{z^L-1} &= \begin{cases}
      \frac{C_0}{z-1} + \sum_{k=1}^{\lfloor L/2 \rfloor} \left( \frac{C_k}{z-\omega^k} + \frac{\overline{C_k}}{z-\overline{\omega^k}} \right), & 2 \!\not|\, L \\[2ex]
      \frac{C_0}{z-1} + \frac{C_{L/2}}{z+1} + \sum_{k=1}^{L/2-1} \left( \frac{C_k}{z-\omega^k} + \frac{\overline{C_k}}{z-\overline{\omega^k}} \right), & 2 \,|\, L
    \end{cases} \\[2ex]
    &= \begin{cases}
      \frac{C_0}{z-1} + \sum_{k=1}^{\lfloor L/2 \rfloor} \frac{2C_kz-2\Re(C_k\overline{\omega^k})}{z^2-2\Re(\omega^k)z+1}, & 2 \!\not|\, L \\[2ex]
      \frac{C_0}{z-1} + \frac{C_{L/2}}{z+1} + \sum_{k=1}^{L/2-1} \frac{2C_kz-2\Re(C_k\overline{\omega^k})}{z^2-2\Re(\omega^k)z+1}, & 2 \,|\, L
    \end{cases}
  \end{align*}
  Let us take a closer look at the terms under the summation signs.
  Since $C_k \in \R$ we have $\Re(C_k\overline{\omega^k})/C_k = \Re(\omega^k)$
  so
  \begin{align*}
    \frac{2C_kz-2\Re(C_k\overline{\omega^k})}{z^2-2\Re(\omega^k)z+1}
      &= C_k \cdot \frac{2z-2\Re(\omega^k)}{z^2 -2\Re(\omega^k)z +1} \\
    C_k \int \frac{\left( z^2 -2\Re(\omega^k)z +1 \right)'}{z^2-2\Re(\omega^k)z+1} dz
      &= C_k \cdot \ln \left( z^2-2\Re(\omega^k)z+1 \right)
  \end{align*}
  Integrating term by term we get
  \begin{align*}
    \int \frac{S_{\sigma_i, \sigma_j}'(z)}{S_{\sigma_i, \sigma_j}(z)}
      &= \int \frac{-P(z)}{z^L-1} \\
    &= \begin{cases}
      C_0\ln(z-1) + \sum_{k=1}^{\lfloor L/2 \rfloor} C_k \ln \left( z^2-2\Re(\omega^k)z+1 \right)  , & 2 \!\not|\, L \\[2ex]
       C_0\ln(z-1) + C_{L/2}\ln(z+1) + \sum_{k=1}^{L/2-1} C_k \ln \left( z^2-2\Re(\omega^k)z+1 \right) , & 2 \,|\, L
    \end{cases}
  \end{align*}
  Taking exponent and setting $A_k := C_k$ provides us with the result.
\end{proof}

\end{document}
